\section{Methods}

\subsection{The estimator}
\label{sec:estimator}

The estimator, \populacefit{}, models $P(\text{targets} \mid \text{predictors})$ with a
quantile-regression-forest estimator \citep{meinshausen2006quantile,
johnson2024quantile} wrapped in three mechanisms.

\paragraph{Weighted bootstrap.} A quantile regression forest predicts by
lookup. Each tree is grown until every leaf holds one or a few training
rows, and the forest predicts for a new row from the target values of the
training rows that share its leaves. Passing per-record weights to the
forest when it is fitted can change where the trees
split, but each stored training row still counts once when a value is read
out, so the predicted distribution is not weighted. The estimator therefore
puts the weights into the data. Before each forest is grown, training rows
are resampled with replacement with probability proportional to weight, so a
household that stands for many others appears many times in the leaves.
Draws from the fitted forest then follow the weighted conditional
distribution. Resampling to undo an
unequal-probability design is an established device
\citep{hinkins1997inverse, dong2014nonparametric}; Section~\ref{sec:background-weights}
places this one in that literature.

\paragraph{Sign gates.} Economic variables are rarely smooth. Most adults
hold zero dividend income; net worth is negative for some households, exactly
zero for others, and spans several orders of magnitude when positive. A
single regressor fit across such a target either drops a tail or interpolates
across the empty gap at zero. The classical remedy is a two-part model, also
called a hurdle model: one model for whether the value is zero and another
for its size when it is not
\citep{mullahy1986specification, lambert1992zero}. The estimator
generalizes the two-part idea to the target's \emph{sign support}, meaning
whichever of $\{\text{negative}, \text{zero}, \text{positive}\}$ occur in
the data, and imputes in four steps:
\begin{enumerate}
\item \emph{Detect the sign classes.} Read which sign classes occur in the
training data. Detection is deliberately unweighted: which signs
exist is a fact about the variable, not about the sampling design,
and a rare but real class (negative net worth, say) must not vanish because
its weight share is small.
\item \emph{Fit the gate.} A classifier learns the probability of each sign
class given the predictors. It is a histogram gradient-boosting classifier
\citep{pedregosa2011scikit} at its default settings, and the version
evaluated here offers no other. The gate takes the weights directly, as
per-record weights the classifier honors, and not by resampling the training
rows, because resampling could delete a rare sign class outright, the
failure step 1 exists to prevent.
\item \emph{Fit magnitudes within each sign.} A separate quantile-regression
forest models the magnitude of each nonzero sign class, so no forest ever
straddles the gap at zero.
\item \emph{Draw.} For each receiver row, sample a sign class from the
gate's predicted probabilities. Taking the modal class would collapse the
zero mass. If the sampled class is nonzero, draw a magnitude from that
class's forest.
\end{enumerate}
For dividend income the sign classes are $\{0, +\}$ and the gate is a zero
hurdle; for net worth they are $\{-, 0, +\}$ and the gate is three-way.
The gate's probabilities are those of the classifier as fitted on the
training rows: the only behavior at the version evaluated here, and still
the package's default. For a rare sign class such a gate can be
overconfident on rows it has not seen, and later versions of \populacefit{}
add an opt-in out-of-fold calibration of the gate for that case
\citep{microcosm2026}. The tasks in this paper contain no rare sign class on
which the methods differ, so the gate's behavior on rare classes is not
evaluated here.
Section~\ref{sec:results-cps} compares the gated estimator with ungated
baselines on zero-inflated targets.

\paragraph{Sequential chaining.} Targets are imputed in order, each
conditioning on the predictors plus the targets already drawn
\citep{raghunathan2001multivariate, vanbuuren2011mice}, so dependence among
imputed variables survives.

\paragraph{Interface contract.} Fitting over a \populace{} frame reads the
frame's typed survey weights by default. Fitting over a plain data frame,
the standalone path, requires the weights explicitly (a weight column name,
a weight vector, or the literal \texttt{"none"}), and omitting them is an
error, not a silent unweighted fit. Under this contract the weight-blind fit
of Section~2.3 cannot happen by omission.

\subsection{Ablations and baselines}
\label{sec:ablations}

\landscapefloat{table}{The methods compared}{tab:methods}{%
\fitwidth{\tablefont\setlength{\tabcolsep}{3pt}%
\begin{tabular}{lllll}
\hline
Method & Weights & Gates & Chaining & Implementation \\
\hline
\populacefit{} & design & yes & yes & \populacefit{} \\
\quad -- unweighted & none & yes & yes & \texttt{weights="none"} \\
\quad -- unchained & design & yes & no & targets fit independently \\
\quad -- ungated & design & no & yes & single-sign path (Section~\ref{sec:ablations}) \\
\quad -- ungated, unchained & design & no & no & both of the above \\
\quad -- net worth derived & design & yes & yes & net worth from three imputed components \\
\quad -- unchained, net worth derived & design & yes & no & the same, components fit independently \\
QRF (\microimpute{}) & design & no & yes & \microimpute{} QRF, as shipped \\
OLS & design & no & no & \microimpute{} \\
Linear quantile regression & design & no & no & \microimpute{} \\
NND hot deck & standardization only & \multicolumn{2}{l}{via donation} & \pystatmatch{}, nearest donor \\
Random hot deck & standardization only & \multicolumn{2}{l}{via donation} & \pystatmatch{}, RANDwNND, equal chances \\
Weighted random hot deck & standardization, donor choice & \multicolumn{2}{l}{via donation} & \pystatmatch{}, weighted draw (RANDwNND) \\
Constrained hot deck & standardization, exact ties & \multicolumn{2}{l}{via donation} & \pystatmatch{}, capped donor use \\
Weighted marginal draw & design & --- & --- & unconditional weighted draw \\
\hline
\end{tabular}}}{%
Notes: Every fit that has a weights column uses the survey's design weights;
the one exception is the unweighted ablation, whose sole difference from the
full estimator is \texttt{weights="none"}. The four ablations are the full
estimator with one or two design choices switched off and nothing else
changed. The two net-worth-derived variants run only in the production-scale
profile, where the targets form a balance sheet. The \microimpute{} forest is
an external baseline run as shipped: it applies the design weights by one
resample of its training rows and chains its targets, but it has no gate
and it draws through its own sampler, so it differs from the full
estimator in more than one respect and is not one of its ablations. The hot decks donate all target
values from one matched donor record, which preserves zero masses and
cross-target structure by construction rather than by modeling them; they
differ in which donor is taken, and only the weighted random hot deck draws
its donor by weight. The random hot deck is its control: the same draw among
the nearest donors with the weights left out. The constrained hot deck uses
the weights only to break exact ties in distance.}

Table~\ref{tab:methods} lists the methods. Each design choice is
evaluated by switching it off with everything else held fixed: the full
estimator versus itself with \texttt{weights="none"} (the weighted
bootstrap's contribution), versus fitting each target independently
(chaining's contribution), and versus itself with the gates off (gating's
contribution).

The package has no option that switches its gates off, but it builds a
gate only for a target whose training values fall in more than one sign
class; a target whose values are all positive is fitted with one forest and
no gate. The gates-off ablation uses that path. For each target with more
than one sign class it adds a constant to every training value, chosen so
that the smallest value becomes 1, fits the estimator as usual, and subtracts
the constant from the drawn values. A regression tree chooses its splits
from differences between target values, which adding a constant does not
change, so the result is the estimator without its gates. The weighted
bootstrap, the forests, the draw, and the chaining are unchanged, and a
single forest covers all rows, zeros included, where the full estimator has
a classifier and a forest per sign. That equivalence is exact in arithmetic
and approximate on a computer: with a large constant, rounding can change
which of two nearly equal splits a tree takes, which perturbs individual
trees the way a different random seed would. A target that already has one
sign class is not shifted, and there the ablation is the full estimator.
Switching off gates and chaining together completes a two-by-two design,
which shows where the two interact.

Where the targets satisfy an accounting identity, a practitioner would impute
all but one of them and compute the last. The net-worth-derived variants do
so in the production-scale profile: they impute the three balance-sheet
components, chained or independently, and set net worth to their signed sum.
They separate what chaining contributes through the identity from what it
contributes through the dependence among the components.

Baselines are implementations of standard methods:
quantile-regression-forest, ordinary-least-squares, and
linear quantile-regression imputers from \microimpute{}
\citep{policyengine2025microimpute}, and distance hot-deck
statistical matching from \pystatmatch{}, a Python port of the R package StatMatch
\citep{dorazio2025statmatch}. The three \microimpute{} imputers return
conditional quantiles, and the benchmark turns each into a sampler in the
same way. The model predicts a grid of 99 quantile levels (0.01 to 0.99),
and each receiver row takes the prediction at one level drawn uniformly
from that grid. The level is drawn once per row and applied to every target
of the row. The \microimpute{} forest is built on the same
forest library as the full estimator \citep{johnson2024quantile}, but it is
not one of its ablations. It applies the design weights by resampling its
training rows once, at a fixed seed, in proportion to weight. It chains its
targets. And it turns each requested quantile level into a draw from a
ten-level conditional quantile grid (levels $1/11$ to $10/11$) at a randomly
chosen index. Under the level sampling just described, that draws
the lowest and highest levels 25 and 21 percent of the time, against 10
percent each for a uniform draw. Its generator is re-seeded identically for
every target, so all of a row's chained targets are read at the same index.

The four hot decks share one distance: Euclidean on the matching variables,
which \pystatmatch{} standardizes by weighted moments when the donor weights
are unequal. They differ in the donor taken. The nearest-neighbour hot deck
takes the closest donor, so the weights do not enter its choice. The weighted
random hot deck draws one of the $\lceil\sqrt{n}\rceil$ nearest of the $n$
donors with probability proportional to its weight
\citep{dorazio2006statistical}. The random hot deck draws among the same
donors with equal probability. It is a control and not a recommendation:
the pair separates what the weights contribute from what the draw itself
does, and under equal weights the two coincide. The constrained hot deck finds the assignment
of least total distance in which no donor serves more than
$\lceil n_{\text{receivers}} / n \rceil$ receivers. The cap counts receivers,
and the weights enter that assignment only to break exact ties in distance,
in favour of the higher-weight donor. Its assignment problem is dense, and at the
population view's 20,000 receivers one split took thirteen minutes against
seconds for every other method, so it is run on the within-survey tasks only.

\subsection{Joint metrics: the population-view harness}
\label{sec:harness}

Per-variable metrics cannot settle the question this paper asks, because a
marginal metric cannot tell a conditional method from an unconditional one.
For the joint distribution we use the population-view harness of
\citet{ghenis2026popdgp}, of which this paper is the first application. This
section summarizes what is needed to read the results; the framework's
general treatment is in that work.

\landscapefloat{figure}{The population-view harness}{fig:harness}{%
\fitwidth{\tablefont% The population-view harness, schematically. Hand-authored (a design
% description, not a results exhibit); pdflatex-compatible TikZ only.
% All text is 10 pt sans serif (IJM: text in figures in Arial or similar,
% size 10).
\begin{tikzpicture}[
    font=\sffamily\fontsize{10}{12}\selectfont,
    node distance=0.9cm and 0.9cm,
    box/.style={draw, rounded corners=1.5pt, align=center,
                inner sep=4pt, minimum height=2.4em},
    faint/.style={draw=black!60, text=black!70},
    flow/.style={-{Stealth[length=2mm]}, thick},
    faintflow/.style={-{Stealth[length=2mm]}, thick, black!55, dashed},
    lab/.style={font=\sffamily\fontsize{10}{12}\selectfont\itshape, align=center},
]

% Latent population
\node[box, minimum width=4.2cm] (pop)
    {latent population $P$\\[-1pt]joint of demographics,
     income, wealth, \ldots};

% Two survey views
\node[box, below left=1.35cm and -2.9cm of pop, minimum width=4.2cm] (scf)
    {\scf{} view $V_{\text{\scf{}}}(P)$\\[-1pt]demographics
     $+$ wealth; list design, \scf{} weights};
\node[box, below right=1.35cm and -2.9cm of pop, minimum width=4.2cm] (cps)
    {\cps{} view $V_{\text{\cps{}}}(P)$\\[-1pt]demographics
     $+$ income; area design, \cps{} weights};

\draw[flow] (pop) -- (scf) node[midway, left=14pt, lab] {observed\\through};
\draw[flow] (pop) -- (cps) node[midway, right=14pt, lab] {observed\\through};

% SCF splits
\node[box, below left=1.1cm and -3.6cm of scf, minimum width=2.5cm] (donor)
    {donor split};
\node[box, below right=1.1cm and -2.6cm of scf, minimum width=2.5cm] (holdout)
    {holdout $H_{\text{\scf{}}}$\\[-1pt]never used upstream};

\draw[flow] (scf) -- (donor) node[midway, left=8pt, lab] {80\%};
\draw[flow] (scf) -- (holdout) node[midway, right=8pt, lab] {20\%};

% Candidate assembly
\node[box, below=2.4cm of cps, minimum width=4.2cm] (cand)
    {candidate file $Q$\\[-1pt]real \cps{} demographics
     $+$ imputed wealth,\\[-2pt]under \cps{} weights};

\draw[flow] (cps) -- (cand);
\draw[flow] (donor.south) to[out=-90, in=180]
    node[pos=0.55, below=3pt, lab] {fit wealth conditionals}
    (cand.west);

% Scorecard
\node[box, below=3.2cm of donor, anchor=north west, xshift=-0.3cm,
      minimum width=9.6cm] (score) at (donor.south -| donor.west)
    {score $V_{\text{\scf{}}}(Q)$ against $H_{\text{\scf{}}}$, in the
     \scf{} view's own variables\\[1pt]
     energy distance (joint) \;$\cdot$\; coverage (support)
     \;$\cdot$\; classifier AUC (omnibus)\\[1pt]
     tail block (uncapped $W_1$/sd, q90/q99 ratios)};

\draw[flow] (cand.south) to[out=-90, in=0]
    node[pos=0.5, right=4pt, lab] {projected\\through $V_{\text{\scf{}}}$}
    (score.east);
\draw[flow] (holdout) -- (holdout |- score.north);

% Sampling-noise floor
\draw[faintflow] (donor) -- (donor |- score.north)
    node[midway, left=2pt, lab, text=black!55]
    {donor split\\scored as a\\reference candidate\\(sampling-noise floor)};

\end{tikzpicture}
}}{%
Notes: Two surveys observe the same population. The Survey of Consumer
Finances (\scf{}) records demographics and wealth, under a design that
oversamples wealthy households. The Current Population Survey (\cps{})
records demographics and income, under an area-probability design. A
candidate file joins real \cps{} demographics to imputed wealth. The harness
restricts it to the variables the \scf{} observes and scores it against a
held-out \scf{} sample on four metrics. The held-out sample is never used to
fit a method, and the donor split itself is scored as a reference to show
the size of sampling noise.}

The harness treats each survey as a \emph{view} of one population: a subset
of its variables, collected under one sampling design. A candidate file is
scored by restricting it to a view's variables and comparing it, under its
weights, with a held-out sample of that survey, the \emph{holdout}. Here
there is one view, the \scf{}'s. Candidates are \cps{} households with
imputed wealth, and they are compared with held-out \scf{} households on the
shared predictors and wealth together (Figure~\ref{fig:harness}). The
holdout is never used to fit a method, so no candidate is scored against
data it was built from.

The harness reports four metrics. Two of them compare every pair of points,
which is costly on large samples. For those two, each sample larger than
2,048 points is first reduced to 2,048 by a resample drawn in proportion to
weight. We call this limit the \emph{subsample cap}, and we say a metric is
capped when it applies and uncapped when the metric uses every row. Write
$X = \{(x_i, a_i)\}$ for the candidate sample and $Y = \{(y_j, b_j)\}$ for
the holdout, with weights normalized to sum to one.

\paragraph{Energy distance (capped; lower is better, 0 is perfect).}
The weighted energy distance \citep{szekely2013energy}
\begin{equation}
\mathcal{E}(X, Y) = 2 \sum_{i,j} a_i b_j \lVert x_i - y_j \rVert
- \sum_{i,i'} a_i a_{i'} \lVert x_i - x_{i'} \rVert
- \sum_{j,j'} b_j b_{j'} \lVert y_j - y_{j'} \rVert
\label{eq:energy}
\end{equation}
is the sample form of a strictly proper scoring rule
\citep{gneiting2007strictly}: in expectation it is smallest only for the
true distribution.

\paragraph{Coverage (capped; higher is better, 1 is perfect).} \prdc{}
coverage \citep{naeem2020reliable} is the weighted fraction of holdout
points with at least one candidate point inside their
$k$-nearest-neighbour radius. It asks whether the candidate reaches every
neighbourhood the holdout occupies. It depends on the candidate only through
which points it contains, so in principle reweighting the candidate cannot
change it. As implemented, the weight-proportional resample to the cap makes
that hold only approximately.

\paragraph{Classifier two-sample test (uncapped; 0.5 is indistinguishable).}
A gradient-boosted classifier is trained to tell candidate rows from
holdout rows at equal class mass, and we report its cross-validated area
under the ROC curve (AUC) \citep{lopezpaz2017revisiting}. We abbreviate the
test as C2ST. An AUC of 0.5 means the classifier cannot separate the files;
1.0 means it always can, so a lower AUC is better. The classifier can pick
up any difference between the two files that it is able to learn, including
broken dependence between variables whose marginals are each right, so it
serves as a single overall test.

\paragraph{Tail metrics (uncapped; a ratio of 1 is perfect).} For each
imputed column the harness reports its weighted Wasserstein-1 distance
(Equation~\ref{eq:w1}) to the holdout, scaled by the holdout's weighted
standard deviation, and the ratios of its weighted 90th and 99th percentiles
to the holdout's. These were added after the first three metrics failed to
capture an imputed 99th percentile of wealth that was double the holdout's
(Section~\ref{sec:results-harness}). Wealth and income have heavy right
tails, where a small share of households holds much of the total.

We call the first three the \emph{joint metrics}, because each compares the
shared predictors and wealth together, and the fourth the tail metrics,
which look at one imputed column at a time. A \emph{sampling floor}, the
donor split itself scored as a candidate, shows what ``as good as another
sample of the survey'' means on each metric, including how much a tail ratio
moves under sampling noise alone.

\subsection{Within-survey metrics}
\label{sec:metrics}

All metrics are computed on held-out receiver rows under survey weights,
with one exception noted under pinball loss.

\paragraph{Pinball loss.} Pinball loss (also called quantile loss; we use
\emph{pinball} throughout) scores predicted quantiles against realized
values \citep{koenker2005quantile}.
For a predicted quantile $\hat{q}_\tau$ and realized value $y$, the
loss is $\rho_\tau(y - \hat{q}_\tau)$ with
$\rho_\tau(u) = u\,(\tau - \mathbf{1}[u < 0])$; we average it, weighted,
over the decile grid $\tau \in \{0.1, \ldots, 0.9\}$. A single stochastic
draw carries no conditional quantiles, so $\hat{q}_\tau$ here is the
$\tau$-quantile of the method's pooled receiver draw, read with every
receiver row counted once, and it is scored against every held-out value
under the receiver weights. As computed, pinball loss is therefore a
marginal score. It does not test conditional calibration and cannot
separate a conditional method from an unconditional one. Because the
predicted quantiles are read without weights, it also favors a draw whose
unweighted quantiles already match the weighted population.

\paragraph{Wasserstein-1 distance.} Marginal fit is scored by the weighted
Wasserstein-1 distance to the weighted donor distribution,
\begin{equation}
W_1(F, G) = \int_0^1 \bigl| F^{-1}(u) - G^{-1}(u) \bigr|\, du ,
\label{eq:w1}
\end{equation}
the average absolute gap between the two distributions' quantile functions,
or equivalently the area between their CDFs. The subscript is the order:
$W_1$ averages absolute quantile gaps in the variable's own units (dollars
here), where $W_2$ would square them. We use order 1 because it is
interpretable in dollars and does not let a single extreme quantile
dominate the score quadratically.

\paragraph{Zero-share error.} The absolute difference between the
candidate's and the donor's weighted share of exact zeros. It measures how
well a method preserves the zero mass of a zero-inflated target.

\paragraph{Reweighting fragility.} Fragility measures how much of a total
one record could come to carry after a later reweighting, which is the
landmine problem of Section~1. Suppose a calibration multiplies each
record's weight $w_i$ by a factor between $1/\kappa$ and $\kappa$. The
reweighting that concentrates the total most multiplies the largest
contributor's weight by $\kappa$ and divides every other weight by $\kappa$.
That record's share of the total is then
$\kappa^2 c^* / (\kappa^2 c^* + S - c^*)$, where $c^*$ is the largest
contribution $w_i |a_i|$ of a record with value $a_i$, and $S$ is the sum of
all contributions. We report this share at $\kappa = 5$ for every method,
next to the same quantity for the held-out true values. The truth's value is
the exposure a faithful method should match, not minimize: a much lower
value means the method spreads the total too evenly, and a much higher one
means it has created an outsized record.

\paragraph{Paired classifier test.} The metrics above look at one imputed
variable at a time. They cannot tell a method that follows the predictors
from one that draws each variable from the right distribution and ignores
them. Within a survey the receivers' true values are known, so conditional
fit can be tested directly.

We make two copies of the receiver rows. One copy carries each row's true
target values and the other carries the imputed ones. The predictors and the
weights are the same in both. A gradient-boosted classifier is trained to
tell the two copies apart from the predictors and the targets together, and
we report its cross-validated AUC under the receiver weights. Both copies of
a row sit in the same fold, so the classifier is never tested on a row whose
other copy it was trained on. Because the predictors are identical in the
two copies, the classifier can separate them only through how the targets
vary with the predictors and with each other. An AUC of 0.5 means it cannot
separate them, and a lower AUC is better.

This test differs from the classifier two-sample test of
Section~\ref{sec:harness}, which compares two different sets of households.
To show how far its AUC moves by chance, we repeat it with the labels of a
random half of the rows swapped. If the imputed values followed the true
conditional distribution, either copy of a row could carry either label, so
the repeat is a draw from the test's distribution when the imputation is
right. We call its result the null AUC.

\paragraph{Wasserstein-1 distance within groups.} The classifier counts
households. Every receiver has the same say in it, whatever its wealth, so a
method can score well while misplacing a few very large values, and those
are what move a total. We therefore pair the classifier with a conditional
measure in dollars. The receivers are divided into groups of similar
households. Within each group we compute the weighted $W_1$ distance between
the imputed values and the same rows' true values. We report the average of
the group distances, weighted by each group's total weight.

In the \scf{} task the groups are the fifths of household income crossed
with three age bands (under 40, 40 to 64, and 65 or over), 15 in all. In the
\cps{} task they are the adults with no earnings and the thirds of those with
earnings, crossed with the same age bands, 12 in all. The cut points are
weighted quantiles of the task's full sample, so every split and every
method is scored on the same groups. A group with fewer than five receivers
in a split is left out of that split's average.

\paragraph{Imputed total.} Beside the within-group distance we report each
target's imputed total as a multiple of the survey's. It is the weighted
mean of the imputed values over the receivers, divided by the weighted mean
of the true values in the task's full sample. A value of 1 reproduces the
survey's total.

\subsection{Protocol}
\label{sec:protocol}

Every task uses an 80/20 donor/receiver holdout repeated over ten seeds, with
splits paired across methods: the split is a pure function of the seed, so
every method fits and is scored on identical partitions. Task inputs are
frozen once and pinned by hash; only the method varies within a sweep. All
results tables are generated from the sweep artifacts by the repository's
command line; no result is hand-entered. Where a method is unavailable or a
metric undefined for a task, the cell is reported as absent rather than
omitted.

The paired classifier test, the within-group distance and the imputed total
were added after the 2025 replication had been scored
(Section~\ref{sec:results-conditional}). They refit every method on the same
splits with the same seeds, so they score the same draws as the other
within-survey tables. On the three survey tasks they use 30 seeds, the ten
above and twenty more. On the OpenML datasets they use the ten.
